% Every operator-visible message string, with cause and action.
%
% DERIVED FROM FIRMWARE. Refresh against:
%   Firmware/Core/Src/robot.cpp                 (request names, outcomes, faults, reports)
%   Firmware/Core/Src/user_interface_module.cpp (menu feedback)
%   Firmware/Core/Src/menu.cpp                  (ERROR/WARNING headers, severity behaviour)

% ------------------------------------------------------------- operations --

\subsection*{Operation names}

Operations announce themselves on the top line and report their outcome on the
line below. These six names can appear:

\begin{center}\small
\begin{tabular}{@{}L{32mm}L{20mm}L{70mm}@{}}
\toprule
\textbf{Name} & \textbf{Severity} & \textbf{What it is} \\
\midrule
\msg{HOMING}      & critical & Y-axis reference search, run automatically at power-up \\
\msg{BONDER INIT} & critical & Bond-head initialisation, run automatically after homing \\
\msg{US TEST}     & normal   & Ultrasonic test, started with \keycap{TEST} \\
\msg{FORCE TEST}  & normal   & Force calibration run, started with \keycap{SETUP} \\
\msg{TACH. CAL.}  & normal   & Tachometer calibration, started from the Settings page \\
\msg{BONDING}     & normal   & A bond cycle \\
\bottomrule
\end{tabular}
\end{center}

A \textbf{critical} operation that fails locks the machine. A \textbf{normal}
one that fails only warns; you can retry it.

\subsection*{Operation outcomes}

\begin{center}\small
\begin{tabular}{@{}L{42mm}L{80mm}@{}}
\toprule
\textbf{Message} & \textbf{Meaning} \\
\midrule
\msg{STARTED}             & The operation has begun. \\
\msg{OPERATION SUCCEEDED} & It finished normally. \\
\msg{OPERATION FAILED}    & It stopped early. A fault line above says why. \\
\msg{CAN'T START}         & It could not begin --- normally because the machine is already locked. \\
\msg{ERROR SYSTEM LOCKED} & A critical operation failed. The machine will not move again until it is reset. \\
\bottomrule
\end{tabular}
\end{center}

% ------------------------------------------------------------------ faults --

\subsection*{Faults}

When an operation fails, the reason appears above the outcome line.

\begin{center}\small
\begin{tabular}{@{}L{40mm}L{40mm}L{45mm}@{}}
\toprule
\textbf{Message} & \textbf{Cause} & \textbf{What to do} \\
\midrule

\msg{LOW BONDING POWER} &
The requested ultrasonic energy was not delivered before \menuitem{Bond Timeout}
expired, or the scan could not find resonance. &
Run \keycap{TEST}. If the reported frequency is outside
\menuitem{Scan Start}--\menuitem{Scan Stop}, widen the window. Otherwise check
the transducer and its cable, then raise \menuitem{Power} or lower
\menuitem{Energy}. \\
\addlinespace

\msg{COIL CURR. UNSTABLE} &
The force coil could not reach or hold its commanded force. &
Check the coil wiring and that the bond lever moves freely. Reduce the
commanded force. If it persists at every force, the coil drive needs service. \\
\addlinespace

\msg{Z POS. NOT SETTLED} &
The Z axis did not reach its commanded height, or the position loop dropped
out. &
Check that nothing is obstructing the bond head. Verify \menuitem{Reset Height}
and \menuitem{Overtravel} are inside the Z travel. Run tachometer calibration. \\
\addlinespace

\msg{TIMEOUT} &
The machine waited for an expected event that never arrived --- most often a
touchdown that never happened or a button that was never released. &
Confirm the workpiece is within reach at the search height and that the contact
sensor is connected. Retry the cycle. \\

\bottomrule
\end{tabular}
\end{center}

% ------------------------------------------------------------------ system --

\subsection*{System messages}

\begin{center}\small
\begin{tabular}{@{}L{40mm}L{85mm}@{}}
\toprule
\textbf{Message} & \textbf{Meaning and what to do} \\
\midrule
\msg{BUSY! TRY LATER} &
You pressed \keycap{SETUP}, \keycap{TEST} or \keycap{CLAMP OPEN}, or started
tachometer calibration, while the machine was already running something. Wait
for the current operation to finish and press again. \\
\bottomrule
\end{tabular}
\end{center}

\subsection*{Ultrasonic test report}

\keycap{TEST} reports three lines, shown for ten seconds:

\begin{center}\small
\begin{tabular}{@{}L{34mm}L{88mm}@{}}
\toprule
\textbf{Line} & \textbf{Meaning} \\
\midrule
\msg{Freq:58.123 kHz} & Series resonance found by the sweep. Must fall between
                        \menuitem{Scan Start} and \menuitem{Scan Stop}. \\
\msg{Q:12.3}          & Quality factor of the transducer. Watch this over time;
                        a falling Q signals a loosening horn or a failing stack. \\
\msg{Power:0.5 W}     & Power actually delivered during the test. \\
\bottomrule
\end{tabular}
\end{center}

% -------------------------------------------------------------------- menu --

\subsection*{Menu messages}

These confirm or reject something you did in the menus. All clear on their own
or on the next key press.

\begin{center}\small
\begin{tabular}{@{}L{42mm}L{80mm}@{}}
\toprule
\textbf{Message} & \textbf{Meaning} \\
\midrule
\msg{Configuration saved}   & The active configuration was written to memory. \\
\msg{Configuration added}   & A new configuration was created. \\
\msg{Settings saved}        & The Settings page was written to memory. \\
\msg{Force offset saved}    & Your gauge reading was stored as the force correction. \\
\msg{Save failed}           & The configuration could not be written. Storage may be full --- 32 configurations is the limit. \\
\msg{Load failed}           & The configuration could not be read. \\
\msg{Add failed}            & The configuration could not be created. Storage is likely full. \\
\msg{Delete failed}         & The configuration could not be removed. \\
\msg{Cannot delete DEFAULT} & \menuitem{DEFAULT} is permanent. Delete one of your own configurations instead. \\
\msg{Not used by protocol}  & You pressed a hotkey for a parameter the active bonding mode does not use. \\
\bottomrule
\end{tabular}
\end{center}

\derivedfrom{\texttt{robot.cpp}, \texttt{user\_interface\_module.cpp},
             \texttt{menu.cpp}}
